% Generated by tools/build_new_dists_anthology_sections.py; do not edit by hand.

\section{Lattice automorphism groups}\label{proof:new-dist:lattices}

\resultbox{\textbf{Canonical testing artifacts.}\par\smallskip Exact backend: \path{new_dists/lattice_sigma_mgfs/verification.py}.\par Regression: \path{new_dists/lattice_sigma_mgfs/test_verification.py}.\par Marimo: \path{marimo_notebooks/lattices.py}.}

\subsection*{Scope and outcome}

For the canonical complex representation of the full automorphism group of a
positive-definite integral lattice, we prove the exact frame-shape
$\sigma$-MGF and its coefficient interpretation by Reynolds projection.  We
then prove the cycle formulas and coefficientwise limits for type-$A$ root
lattices, signed-permutation lattices, and repeated orthogonal sums; give the
finite class reductions for Niemeier, Coxeter--Todd, extremal, and modular
lattices; exhibit explicit rank-$2,3,4$ lattices with automorphism group
exactly $\{\pm I\}$; and prove the contrasting exponential moment growth for
tensor powers.  A reproducible marimo notebook constructs all tractable
matrix groups and compares direct matrix averages, recovered frame shapes,
family cycle formulas, and independently computed common fixed spaces.

\subsection{Canonical representation and the frame-shape theorem}

Let $L$ be a positive-definite integral lattice of rank $d$, put
\[
 G_L=\Aut(L),\qquad V_L=L\otimes_{\Z}\C,
\]
and let $G_L$ act through its canonical lattice matrices.  Positivity makes
$G_L$ finite.  If $\tau=(\tau_1,\ldots,\tau_s)$ is a partition, define
\[
 W_\tau(V_L)=\bigotimes_{j=1}^s\Sym^{\tau_j}V_L,
 \qquad a_\tau(L)=\dim W_\tau(V_L)^{G_L}.
\]
The symmetric series $\MGF_L=\sum_\tau a_\tau(L)m_\tau$ is the lattice
$\sigma$-MGF.

Every $g\in G_L$ has finite order.  Its characteristic polynomial can be
written uniquely in frame-shape form
\begin{equation}\label{nd:lattices:eq:frame}
 \det(xI-g)=\prod_{m\ge1}(x^m-1)^{a_m(g)},
 \qquad \sum_m m a_m(g)=d.
\end{equation}
The exponents are integers and may be negative; cyclotomic cancellation makes
the product a polynomial.  Equivalently,
\begin{equation}\label{nd:lattices:eq:frame-det}
 \det(I-tg)=\prod_{m\ge1}(1-t^m)^{a_m(g)}.
\end{equation}

\begin{theorem}[frame-shape $\sigma$-MGF]\label{nd:lattices:thm:frame}
For every positive-definite integral lattice,
\begin{align}
 \MGF_L(\mathbf t)
 &=\frac1{|G_L|}\sum_{g\in G_L}
   \prod_{i\ge1}\prod_{m\ge1}(1-t_i^m)^{-a_m(g)},
   \label{nd:lattices:eq:element-frame}\\
 &=\sum_{[g]\subseteq G_L}\frac1{|C_{G_L}(g)|}
   \prod_{i,m}(1-t_i^m)^{-a_m(g)},\label{nd:lattices:eq:class-frame}\\
 &=\sum_{[g]}\frac1{|C_{G_L}(g)|}
   \exp\!\left(\sum_{r\ge1}\frac{p_r}{r}\Tr(g^r\mid V_L)\right).
   \label{nd:lattices:eq:trace-frame}
\end{align}
Moreover,
\begin{equation}\label{nd:lattices:eq:coefficient}
 \boxed{[m_\tau]\MGF_L=\dim W_\tau(V_L)^{G_L}.}
\end{equation}
\end{theorem}

\begin{proof}
For any complex matrix $A$,
\begin{equation}\label{nd:lattices:eq:complete}
 \sum_{q\ge0}\Tr(\Sym^qA)t^q
 =\det(I-tA)^{-1}
 =\exp\!\left(\sum_{r\ge1}\frac{t^r}{r}\Tr(A^r)\right).
\end{equation}
Multiply over the variables $t_i$ and average over $G_L$.  The coefficient of
$t_1^{\tau_1}\cdots t_s^{\tau_s}$ is the trace on $W_\tau(V_L)$ of the
Reynolds idempotent
\[
 R_\tau=\frac1{|G_L|}\sum_{g\in G_L}\rho_{W_\tau}(g).
\]
Its image is exactly $W_\tau(V_L)^{G_L}$, so its trace proves
\eqref{nd:lattices:eq:coefficient}.  Equation~\eqref{nd:lattices:eq:frame-det} gives
\eqref{nd:lattices:eq:element-frame}; grouping elements into conjugacy classes gives
\eqref{nd:lattices:eq:class-frame}; and the exponential identity in
\eqref{nd:lattices:eq:complete} gives \eqref{nd:lattices:eq:trace-frame}.
\end{proof}

The frame exponents and power traces determine one another:
\begin{equation}\label{nd:lattices:eq:trace-frame-inversion}
 \Tr(g^r)=\sum_{m\mid r}m a_m(g),\qquad
 a_m(g)=\frac1m\sum_{e\mid m}\mu(m/e)\Tr(g^e).
\end{equation}
Indeed, the sum of the $r$th powers of all $m$th roots of unity is $m$ if
$m\mid r$ and zero otherwise.  Möbius inversion proves the second formula.
This is also the exact recovery algorithm used by the verifier.

\subsection{Universal parity and low-degree diagnostics}

The central automorphism $-I$ belongs to every $G_L$.  It acts on
$W_\tau(V_L)$ by $(-1)^{|\tau|}$; therefore
\begin{equation}\label{nd:lattices:eq:parity}
 \boxed{[m_\tau]\MGF_L=0\quad\text{if }|\tau|\text{ is odd}.}
\end{equation}
Because the lattice form identifies $V_L$ with its dual,
\begin{equation}\label{nd:lattices:eq:low-degree}
 [m_{(1,1)}]\MGF_L=\dim\End_{G_L}(V_L),
 \qquad [m_{(2)}]\MGF_L\ge1.
\end{equation}
If $V_L$ is complex irreducible, Schur's lemma and uniqueness of the invariant
symmetric form give equality to one in both positions.

The fourth multilinear coefficient is already more sensitive.  If
$V_L\otimes V_L=\bigoplus_\rho b_\rho\rho$, then
\begin{equation}\label{nd:lattices:eq:fourth}
 [m_{(1,1,1,1)}]\MGF_L
 =\dim\End_{G_L}(V_L\otimes V_L)=\sum_\rho b_\rho^2.
\end{equation}
Thus irreducibility in degree one does not control a growing family in degree
four.

More generally, Newton's identity gives the exact moment formula
\begin{equation}\label{nd:lattices:eq:moment}
 [m_\tau]\MGF_L
 =\E_{g\in G_L}\prod_j\left[
 \sum_{\lambda\vdash\tau_j}\frac1{z_\lambda}
 \prod_{r\ge1}\Tr(g^r)^{m_r(\lambda)}\right].
\end{equation}
A coefficient of total degree $D$ depends only on the joint moments of
$\Tr(g),\ldots,\Tr(g^D)$.  This observation is the correct bridge from exact
finite formulas to coefficientwise asymptotics.

\subsection{Type-A root lattices}

For
\[
 A_{n-1}=\{(x_1,\ldots,x_n)\in\Z^n:\textstyle\sum_i x_i=0\},
 \qquad n\ge3,
\]
the full automorphism group is $\{\pm1\}\times S_n$ on the standard
hyperplane representation.  If $\sigma$ has $c_m(\sigma)$ cycles of length
$m$, deleting the constant eigenline from the permutation representation
gives
\begin{equation}\label{nd:lattices:eq:a-det}
 \det(I-t\varepsilon\sigma\mid V_{A_{n-1}})
 =\frac{\prod_m(1-(\varepsilon t)^m)^{c_m(\sigma)}}{1-\varepsilon t}.
\end{equation}
Let $m_m(\lambda)$ denote the multiplicity of $m$ in a partition $\lambda$,
and $z_\lambda=\prod_m m^{m_m(\lambda)}m_m(\lambda)!$.  Averaging
\eqref{nd:lattices:eq:a-det} by cycle type proves
\begin{equation}\label{nd:lattices:eq:a-formula}
 \boxed{
 \MGF_{A_{n-1}}(\mathbf t)=\frac12\sum_{\varepsilon=\pm1}
 \sum_{\lambda\vdash n}\frac1{z_\lambda}
 \prod_i\left[(1-\varepsilon t_i)
 \prod_{m\ge1}\bigl(1-(\varepsilon t_i)^m\bigr)^{-m_m(\lambda)}\right].}
\end{equation}

\subsubsection*{Large-rank limit}
For a uniform permutation, the finite collection of cycle counts converges
with all fixed moments to independent $Z_m\sim\operatorname{Poisson}(1/m)$.
Applying the Poisson probability-generating function to
\eqref{nd:lattices:eq:a-formula} yields, coefficientwise,
\begin{equation}\label{nd:lattices:eq:a-limit}
 \boxed{
 \MGF_{A_\infty}=\frac12\sum_{\varepsilon=\pm1}
 \prod_i(1-\varepsilon t_i)
 \exp\!\left[\sum_{m\ge1}\frac1m\left(
 \prod_i(1-(\varepsilon t_i)^m)^{-1}-1\right)\right].}
\end{equation}
Only finitely many cycle moments enter a bounded-degree coefficient, so the
usual finite-dimensional cycle-count convergence is sufficient.

\subsection{Signed-permutation lattices: types B, C, and D}

The cubic lattice $\Z^n$ has the signed-permutation group
$C_2^n\rtimes S_n$.  The same natural matrix group is the Weyl group in types
$B_n$ and $C_n$, and it is the full automorphism group of $D_n$ for $n\ge5$.
The exceptional triality of $D_4$ must be handled by its own finite class sum.

For a permutation cycle of length $\ell$, let $\delta\in\{\pm1\}$ be the
product of its coordinate signs.  Its determinant is $1-\delta t^\ell$.
After averaging the independent sign product, put
\begin{equation}\label{nd:lattices:eq:b-ell}
 B_\ell(\mathbf t)=\frac12\left[
 \prod_i(1-t_i^\ell)^{-1}+\prod_i(1+t_i^\ell)^{-1}\right].
\end{equation}
The cycle index of $S_n$ now gives
\begin{align}
 \MGF_{B_n}(\mathbf t)
 &=\sum_{\lambda\vdash n}\frac1{z_\lambda}
   \prod_{\ell\ge1}B_\ell(\mathbf t)^{m_\ell(\lambda)},
   \label{nd:lattices:eq:b-formula}\\
 \sum_{n\ge0}u^n\MGF_{B_n}(\mathbf t)
 &=\exp\!\left(\sum_{\ell\ge1}\frac{u^\ell}{\ell}B_\ell(\mathbf t)\right).
 \label{nd:lattices:eq:b-egf}
\end{align}
Since $B_\ell=1+$ terms of positive $\mathbf t$-degree, removing the factor
$(1-u)^{-1}$ and taking coefficients proves the stable limit
\begin{equation}\label{nd:lattices:eq:b-limit}
 \boxed{\MGF_{B_\infty}=\exp\!\left(
 \sum_{\ell\ge1}\frac{B_\ell(\mathbf t)-1}{\ell}\right).}
\end{equation}

\subsection{Repeated orthogonal sums and the stable law}

Fix a lattice $L$ with $G=\Aut(L)$ and let $G\wr S_k$ act on
$V_L^{\oplus k}$ by independent lattice automorphisms and component
permutations.  Define
\[
 A_\ell(\mathbf t)=\MGF_L(t_1^\ell,t_2^\ell,\ldots).
\]

\begin{theorem}[marked compound-Poisson law]\label{nd:lattices:thm:wreath-sum}
\begin{align}
 \sum_{k\ge0}u^k\MGF_{G\wr S_k,V_L^{\oplus k}}
 &=\exp\!\left(\sum_{\ell\ge1}\frac{u^\ell}{\ell}A_\ell(\mathbf t)\right),
 \label{nd:lattices:eq:wreath-sum}\\
 \lim_{k\to\infty}\MGF_{G\wr S_k,V_L^{\oplus k}}
 &=\exp\!\left[\sum_{\ell\ge1}\frac{
 \MGF_L(t_1^\ell,t_2^\ell,\ldots)-1}{\ell}\right]
 \label{nd:lattices:eq:wreath-limit}
\end{align}
coefficientwise in $\mathbf t$.
\end{theorem}

\begin{proof}
For a top permutation cycle $C$ of length $\ell$, multiply the $G$-labels
around $C$ to obtain its cycle product $g_C$.  On the associated block,
\[
 \det(I-tx\mid C)=\det(I-t^\ell g_C\mid V_L).
\]
For uniformly independent labels, $g_C$ is uniform in $G$, independently for
distinct cycles.  The exponential cycle-index formula proves
\eqref{nd:lattices:eq:wreath-sum}.  Next write
\[
 \exp\!\left(\sum_{\ell\ge1}\frac{u^\ell A_\ell}{\ell}\right)
 =\frac1{1-u}\exp\!\left(
 \sum_{\ell\ge1}\frac{u^\ell(A_\ell-1)}{\ell}\right).
\]
In each fixed $\mathbf t$-degree, the second factor is a polynomial in $u$.
Its coefficient sequence therefore stabilizes after multiplication by
$(1-u)^{-1}$, and evaluation at $u=1$ proves \eqref{nd:lattices:eq:wreath-limit}.
\end{proof}

This theorem applies to repeated copies of any fixed root lattice, the
Coxeter--Todd lattice, or any other fixed component.  It is a statement about
the natural wreath subgroup; an orthogonal sum can have additional
automorphisms when accidental isometries occur.

\subsection{Isolated and finite lattice families}

\subsubsection*{Rootful Niemeier lattices}
Let $N$ be rootful with root sublattice $R$, Weyl group
$W(R)\triangleleft\Aut(N)$, and umbral quotient
$G^N=\Aut(N)/W(R)$.  Choose a lift $\widetilde\gamma$ of each
$\gamma\in G^N$.  The exact finite reduction is
\begin{equation}\label{nd:lattices:eq:niemeier}
 \boxed{\MGF_N=\frac1{|G^N|}\sum_{\gamma\in G^N}\frac1{|W(R)|}
 \sum_{w\in W(R)}\prod_i\det(I-t_iw\widetilde\gamma)^{-1}.}
\end{equation}
Changing a lift multiplies it by an element of $W(R)$ and merely permutes the
inner sum, so \eqref{nd:lattices:eq:niemeier} is lift-independent.  If a lift cycles
isomorphic root components with length $\ell$ and total diagram twist
$\delta_C$, that component cycle reduces to
\begin{equation}\label{nd:lattices:eq:niemeier-component}
 \frac1{|W(R_a)|}\sum_{x\in W(R_a)}
 \prod_i\det(I-t_i^\ell x\delta_C)^{-1}.
\end{equation}
Thus all 23 rootful cases are finite Weyl/quotient class calculations.  They
are a rank-$24$ collection, not a rank-asymptotic tower.

\subsubsection*{Coxeter--Todd, extremal, modular, and neighbor constructions}
For the rank-$12$ Coxeter--Todd lattice $K_{12}$, and for every specified
extremal or modular lattice, the exact answer is simply
\begin{equation}\label{nd:lattices:eq:isolated}
 \MGF_L=\sum_{[g]\subseteq\Aut(L)}\frac1{|C(g)|}
 \prod_{i,m}(1-t_i^m)^{-a_m(g)}.
\end{equation}
This is complete once the matrix classes or weighted frame-shape table are
supplied.  Lattice extremality, modularity, strong perfection, a neighbor
relation, or an inclusion $L_n\subset L_{n+1}$ does not by itself impose
compatible automorphism representations.  Consequently no universal
coefficientwise limit follows from those geometric labels alone.

The notebook does not fabricate missing class data for these large isolated
groups.  It verifies the determinant and frame-shape mechanism on exhaustive
small matrix groups; a numerical evaluation of \eqref{nd:lattices:eq:niemeier} or
\eqref{nd:lattices:eq:isolated} requires the corresponding complete external class table.

\subsection{Explicit lattices with only scalar symmetry}

The smallest possible full automorphism group is $\{\pm I\}$.  If $L$ has
rank $d$ and this full group, then
\begin{equation}\label{nd:lattices:eq:scalar-series}
 \boxed{\MGF_L=\frac12\left[
 \prod_i(1-t_i)^{-d}+\prod_i(1+t_i)^{-d}\right].}
\end{equation}
In particular,
\begin{equation}\label{nd:lattices:eq:scalar-low}
 [m_{(1,1)}]=d^2,\qquad [m_{(2)}]=\frac{d(d+1)}2.
\end{equation}
The verifier uses the following concrete positive-definite integral Gram
matrices:
\begin{align*}
 Q_2&=\begin{pmatrix}6&-1\\-1&8\end{pmatrix},\\
 Q_3&=\begin{pmatrix}6&1&2\\1&7&-1\\2&-1&10\end{pmatrix},\\
 Q_4&=\begin{pmatrix}6&1&1&2\\1&7&1&2\\1&1&11&1\\2&2&1&13\end{pmatrix}.
\end{align*}

\begin{proposition}\label{nd:lattices:prop:scalar-grams}
For $d=2,3,4$, the lattice with Gram matrix $Q_d$ has full automorphism group
$\{\pm I_d\}$.
\end{proposition}

\begin{proof}
Each matrix is strictly diagonally dominant with positive diagonal, hence
positive definite.  Its Gershgorin lower bound $\delta_d$ satisfies
$v^TQ_dv\ge\delta_d\|v\|_2^2$.  If $A^TQ_dA=Q_d$ and $v_j$ is column $j$ of
$A$, then $v_j^TQ_dv_j=(Q_d)_{jj}$.  Therefore every coordinate of every
column has absolute value at most
\[
 B_d=\left\lfloor\sqrt{\max_j(Q_d)_{jj}/\delta_d}\right\rfloor.
\]
The values $(\delta_d,B_d)$ are $(5,1),(3,1),(2,2)$.  Exhaustive exact
enumeration of the finite boxes $[-B_d,B_d]^d$, retaining precisely the
column tuples with pairings $v_i^TQ_dv_j=(Q_d)_{ij}$, returns two matrices in
each rank: $I_d$ and $-I_d$.  This column-pairing condition is exactly
$A^TQ_dA=Q_d$, so the enumeration is exhaustive.
\end{proof}

Equations~\eqref{nd:lattices:eq:scalar-low} show immediate degree-two divergence in any
rank-growing tower whose full automorphism groups stay scalar.  This is the
basic obstruction to a universal law for arbitrary lattice families.

\subsection{Tensor products of lattice representations}

Let $G=\Aut(L)$ act on $V=V_L$ with character $\chi$.  The local group $G^k$
acts on $V^{\otimes k}$.  Since
\[
 \Tr\bigl((g_1\otimes\cdots\otimes g_k)^r\bigr)
 =\prod_{j=1}^k\chi(g_j^r),
\]
Theorem~\ref{nd:lattices:thm:frame} gives
\begin{equation}\label{nd:lattices:eq:tensor-local}
 \boxed{\MGF_{G^k,V^{\otimes k}}
 =\frac1{|G|^k}\sum_{g_1,\ldots,g_k}
 \exp\!\left[\sum_{r\ge1}\frac{p_r}{r}
 \prod_{j=1}^k\chi(g_j^r)\right].}
\end{equation}
Reordering tensor factors in the multilinear coefficient gives
\begin{equation}\label{nd:lattices:eq:tensor-multilinear}
 \boxed{[m_{(1^s)}]\MGF_{G^k,V^{\otimes k}}
 =\left(\dim(V^{\otimes s})^G\right)^k.}
\end{equation}
If $V$ is irreducible orthogonal of dimension greater than one, then
$\dim(V^{\otimes2})^G=1$, while
\[
 \dim(V^{\otimes4})^G=\dim\End_G(V\otimes V)\ge2.
\]
The last inequality holds because $V\otimes V$ contains the trivial module
and at least one nontrivial constituent.  Hence
\begin{equation}\label{nd:lattices:eq:tensor-growth}
 [m_{(1,1,1,1)}]\MGF_{G^k,V^{\otimes k}}\ge2^k.
\end{equation}
Tensor towers therefore usually exhibit exponential moment divergence,
unlike the repeated-sum law in Theorem~\ref{nd:lattices:thm:wreath-sum}.

If factor permutations are included, write
$x=(g_1,\ldots,g_k;\pi)\in G\wr S_k$.  For a cycle $C$ of $\pi$, let $a_C$
be the ordered cycle product and $d_C(r)=\gcd(|C|,r)$.  The $r$th power splits
$C$ into $d_C(r)$ cycles, each carrying a conjugate of
$a_C^{r/d_C(r)}$.  Contracting indices around those cycles proves
\begin{equation}\label{nd:lattices:eq:tensor-wreath-trace}
 \boxed{\Tr(x^r\mid V^{\otimes k})
 =\prod_C\chi\!\left(a_C^{r/d_C(r)}\right)^{d_C(r)}.}
\end{equation}
Substitution into \eqref{nd:lattices:eq:trace-frame} is the exact tensor-wreath
$\sigma$-MGF.

\subsection{Exact computational verification}

The companion marimo notebook calls a pure-Python exact verifier.  No
floating-point eigenvalue or numerical determinant enters a pass condition.
The cases are grouped by the formulas they test.

\subsubsection*{Matrix constructions}
\begin{itemize}
\item \textbf{Type $A$.}  For $n=3,4,5$, every $\pm\sigma$ is represented on
the integral basis $e_1-e_n,\ldots,e_{n-1}-e_n$.  The matrices preserve the
Gram matrix with diagonal $2$ and off-diagonal $1$.
\item \textbf{Signed permutations.}  Every signed permutation matrix is
enumerated in dimensions $2,3,4$ and checked against the identity Gram
matrix.  These are structural tests of \eqref{nd:lattices:eq:b-formula}; the dimension-$4$
case is not claimed to be the full $D_4$ automorphism group.
\item \textbf{Scalar lattices.}  All automorphisms of $Q_2,Q_3,Q_4$ are
enumerated by the exhaustive bound in Proposition~\ref{nd:lattices:prop:scalar-grams}.
\item \textbf{Repeated sums.}  All $288$ block-monomial matrices of
$\Aut(A_2)\wr S_2$ on $A_2\perp A_2$ are constructed.
\item \textbf{Tensor products.}  Kronecker matrices construct
$\Aut(A_2)^2$ and $\Aut(A_2)\wr S_2$ on $V_{A_2}^{\otimes2}$.  Their abstract
orders are $144$ and $288$, while the matrix-image orders are $72$ and $144$
because $(-I,-I)$ acts trivially.
\end{itemize}

\subsubsection*{Four independent checks}
For each concrete matrix $A$, the direct route computes
$T_r=\Tr(A^r)$ and then
\begin{equation}\label{nd:lattices:eq:newton-audit}
 h_0(A)=1,\qquad qh_q(A)=\sum_{r=1}^qT_rh_{q-r}(A).
\end{equation}
It averages $\prod_jh_{\tau_j}(A)$, which is the Reynolds trace on
$W_\tau$.  Independently, the verifier recovers $a_m(A)$ by
\eqref{nd:lattices:eq:trace-frame-inversion}, checks the polynomial identity
\[
 \det(I-tA)=\prod_m(1-t^m)^{a_m(A)},
\]
including all negative exponents by cross-multiplication, and expands the
frame product.  A third route uses only the family formulas: ordinary cycles,
signed cycle products, marked wreath cycles, or the tensor trace
\eqref{nd:lattices:eq:tensor-wreath-trace}.

Finally, for $\tau=(1),(2),(1,1),(3)$, the code explicitly constructs the
integer matrices on each $\Sym^{\tau_j}V$, forms their tensor products, and
solves all generator-fixed equations over $\mathbb Q$.  This common-kernel
dimension is compared with the Reynolds trace.  Thus the fixed-space check is
not merely another expansion of the averaged determinant.

\subsubsection*{Representative cases and results}
\begin{center}
\small
\begin{tabular}{@{}llrr@{}}\toprule
family & matrix image & dimension & image order\\\midrule
type $A$ & $\Aut(A_2),\Aut(A_3),\Aut(A_4)$ & $2,3,4$ & $12,48,240$\\
signed & $C_2\wr S_2,C_2\wr S_3,C_2\wr S_4$ & $2,3,4$ & $8,48,384$\\
scalar & $\Aut(Q_2),\Aut(Q_3),\Aut(Q_4)$ & $2,3,4$ & $2,2,2$\\
orthogonal sum & $\Aut(A_2)\wr S_2$ & $4$ & $288$\\
local tensor & $\rho(\Aut(A_2)^2)$ & $4$ & $72$\\
tensor wreath & $\rho(\Aut(A_2)\wr S_2)$ & $4$ & $144$\\\bottomrule
\end{tabular}
\end{center}

Selected degree-four outputs illustrate the different asymptotics:
\begin{center}
\small
\begin{tabular}{@{}lrrrr@{}}\toprule
matrix image & $[m_2]$ & $[m_{11}]$ & $[m_{22}]$ & $[m_{1111}]$\\\midrule
$\Aut(A_2)$ &1&1&2&3\\
$\Aut(A_3)$ &1&1&3&4\\
$\Aut(A_4)$ &1&1&3&4\\
$C_2\wr S_n$, $n=2,3,4$ &1&1&3&4\\
$\Aut(Q_2)$ &3&4&9&16\\
$\Aut(Q_3)$ &6&9&36&81\\
$\Aut(Q_4)$ &10&16&100&256\\
$\Aut(A_2)\wr S_2$ on $V^{\oplus2}$ &1&1&4&6\\
$\rho(\Aut(A_2)^2)$ on $V^{\otimes2}$ &1&1&5&9\\
$\rho(\Aut(A_2)\wr S_2)$ on $V^{\otimes2}$ &1&1&4&6\\\bottomrule
\end{tabular}
\end{center}
All odd total-degree values vanish, as required by \eqref{nd:lattices:eq:parity}.  For
the local tensor case, the value $9=3^2$ at $m_{1111}$ verifies
\eqref{nd:lattices:eq:tensor-multilinear}; adjoining the factor swap takes invariants in
the symmetric square of a three-dimensional space and gives
$\binom{3+1}{2}=6$.

The degree-four audit contains twelve partitions per case and three recorded
coefficient routes, for $12\cdot12\cdot3=432$ values.  The fixed-space audit
has $12\cdot4=48$ comparisons.  There are $1{,}250$ distinct matrices; six
frame checks per matrix (one determinant identity and five complete
characters through degree four) give $7{,}500$ elementwise checks.

\subsection{Asymptotic classification and scope}

\begin{center}
\small
\begin{tabularx}{\linewidth}{@{}>{\raggedright\arraybackslash}p{0.25\linewidth}
>{\raggedright\arraybackslash}p{0.29\linewidth}X@{}}\toprule
family & exact formula & asymptotic conclusion\\\midrule
$A_{n-1}$ & \eqref{nd:lattices:eq:a-formula} & Poisson-cycle limit \eqref{nd:lattices:eq:a-limit}\\
$\Z^n$, types $B/C$, $D_n$ for $n\ge5$ & \eqref{nd:lattices:eq:b-formula} & stable signed-cycle law \eqref{nd:lattices:eq:b-limit}\\
$R^{\perp k}$ with component permutations & \eqref{nd:lattices:eq:wreath-sum} & compound-Poisson limit \eqref{nd:lattices:eq:wreath-limit}\\
rootful Niemeier & \eqref{nd:lattices:eq:niemeier} & finite rank-$24$ collection\\
$K_{12}$ and isolated extremal/modular lattices & \eqref{nd:lattices:eq:isolated} & sequence-dependent; repeated sums use \eqref{nd:lattices:eq:wreath-limit}\\
$\Aut(L)=\{\pm I\}$ & \eqref{nd:lattices:eq:scalar-series} & degree-two divergence\\
$V_L^{\otimes k}$ under $G^k$ & \eqref{nd:lattices:eq:tensor-local} & degree-four exponential growth\\
neighbor or laminated towers & \eqref{nd:lattices:eq:class-frame} & requires compatible group data\\\bottomrule
\end{tabularx}
\end{center}

The structural distinction is sharp: orthogonal repeated sums with component
permutations yield stable marked compound-Poisson frame-shape laws, whereas
tensor powers and low-symmetry towers force divergent invariant moments.
The universal object is therefore the weighted frame-shape distribution
\[
 \boxed{\MGF_L=\E_{g\in\Aut(L)}
 \prod_{i,m}(1-t_i^m)^{-a_m(g)},}
\]
together with whatever compatibility a specified lattice tower supplies.

\subsection{Reproducibility}

From the repository root, run
\begin{verbatim}
python3 -m new_dists.lattice_sigma_mgfs.verification
python3 -m pytest -q new_dists/lattice_sigma_mgfs/test_verification.py
marimo edit marimo_notebooks/lattices.py
\end{verbatim}
The LaTeX source is standalone and can be rebuilt in its folder with
\begin{verbatim}
latexmk -pdf -interaction=nonstopmode -halt-on-error \
  lattice_sigma_mgfs_proof_and_verification.tex
\end{verbatim}
No command in this workflow modifies
\texttt{proved\_matrix\_groups/all\_new\_matrix\_groups\_proofs.pdf}.
